Work on the module $\Lambda^{\bullet}(W \oplus R_{2})$ with the occupation basis; the species of the sixteen are those Hypothesis (S) names, through Proposition \claimref{standard-model.under-s-composites-state-space-spinor}: by $(N_{W}, N_{R})$, $e^{c} = (0, 2)$, $Q = (1, 1)$, $u^{c} = (2, 0)$, $d^{c} = (2, 2)$, $L = (3, 1)$, $\nu^{c} = (0, 0)$. Write $s_{i} = v_{i} \wedge w_{i} = n_{i} - \tfrac12$, so that $\omega = s_{1} + s_{2} + s_{3}$, $B - L = 1 - \tfrac23 N_{W} = -\tfrac23\,\omega$ and $T_{3R} = \tfrac12(N_{R} - 1) = \tfrac12(s_{4} + s_{5})$.

(i) $TT^{\dagger}$ is the projector onto the full state and $T^{\dagger}T$ the projector onto the empty state, so $[T, T^{\dagger}] = P_{12345} - P_{\emptyset}$ is diagonal in the occupation basis. In the orthonormal Majorana basis $a_{i} = v_{i} + w_{i}$, $b_{i} = v_{i} - w_{i}$ ($a_{i}^{2} = 1$, $b_{i}^{2} = -1$) the bivectors are trace-orthogonal, and the coefficient of $e_{x}e_{y}$ in the projection of a matrix $M$ is $\mathrm{tr}((e_{x}e_{y})^{-1}M)/32$. A diagonal matrix has trace zero against every bivector that changes some occupation, which leaves only $a_{i}b_{i} = -2s_{i}$; and $\mathrm{tr}(s_{i}(P_{12345} - P_{\emptyset})) = \tfrac12 - (-\tfrac12) = 1$ while $\mathrm{tr}(s_{i}^{2}) = 8$. Hence $\mu(T) = \sum_{i}\tfrac18 s_{i}$, with no mixed pair-rung component. Splitting the sum gives $\mu_{P} = \tfrac18\omega = -\tfrac{3}{16}(B - L)$ and $\mu_{R} = \tfrac18(s_{4} + s_{5}) = \tfrac14 T_{3R}$, and $\tfrac18 X = \tfrac14 T_{3R} - \tfrac{3}{16}(B - L)$ is the same element; $X$ is the thirteenth generator in the form Proposition \claimref{standard-model.majorana-element} records. The split is the one Proposition \claimref{standard-model.automorphisms-preserve-slot-pair-split-are} selects: the automorphisms preserving the slot/pair split are Pati-Salam, whose Lie algebra is $\so(6) \oplus \so(4)$ on the pair and rung bivectors.

(ii) $\mu(T) = \tfrac18(N - \tfrac52)$ with $N$ the total occupation, so $p(\mu(T))$ acts on each occupation state by a scalar depending only on $N$. On the sixteen the Yukawa pairs have $N(e^{c}) = N(Q) = 2$ and $N(L) = N(d^{c}) = 4$. With $\beta$ the symmetric form of $w_{4}$ (the $10$, Proposition \claimref{families.doublet-direction-module-georgi-jarlskog-clebsch}), $B(L, w_{4}p\,e^{c}) = p_{2}\beta(L, e^{c})$, $B(e^{c}, w_{4}p\,L) = p_{4}\beta(e^{c}, L)$, $B(Q, w_{4}p\,d^{c}) = p_{4}\beta(Q, d^{c})$ and $B(d^{c}, w_{4}p\,Q) = p_{2}\beta(d^{c}, Q)$, where $p_{k}$ is the value at $N = k$. The symmetric part is $\tfrac12(p_{2} + p_{4})\beta$ on both pairs, a multiple of the $10$ form.

(iii) $w_{4}\mu_{P} = \tfrac18 w_{4}\omega = -\tfrac{3}{16}w_{4}(B - L)$, which Proposition \claimref{standard-model.grade-three-georgi-jarlskog-direction} shows to be pure grade three with antisymmetric form.

(iv) $w_{4}s_{4} = w_{4}(n_{4} - \tfrac12) = \tfrac12 w_{4}$, since $w_{4}n_{4} = w_{4}$. So $w_{4}\mu_{R}\mu_{P} = \tfrac{1}{64}w_{4}(s_{4} + s_{5})\omega = \tfrac{1}{64}\bigl(\tfrac12 w_{4}\omega + w_{4}s_{5}\omega\bigr)$. By Proposition \claimref{families.doublet-direction-module-georgi-jarlskog-clebsch}(i), $T_{\mathrm{GJ}} = \tfrac43 w_{4}s_{5}\omega$, so the composite is $\tfrac{1}{64}\bigl(\tfrac34 T_{\mathrm{GJ}} - \tfrac34 w_{4}(B - L)\bigr) = \tfrac{3}{256}\bigl(T_{\mathrm{GJ}} - w_{4}(B - L)\bigr) = \tfrac{3}{256}w_{4}(B - L)(\chi - 1)$. $T_{\mathrm{GJ}}$ has symmetric form with ratio $3$ (that proposition, (ii)) and $w_{4}(B - L)$ antisymmetric form (Proposition \claimref{standard-model.grade-three-georgi-jarlskog-direction}); their entries on $(L, e^{c})$ and $(Q, d^{c})$ have equal moduli, $1$ and $\tfrac13$. Since $\chi - 1 = -2$ on $N_{R}$ even and $0$ on $N_{R} = 1$, the composite annihilates the doublet sector.

The script builds the module as exact integer matrices with the pairing of Lemma \claimref{standard-model.invariant-pairing-five-modes}, checks the pairing's $\so(10)$ invariance on sample bivectors, reproduces the forms of $w_{4}$, $T_{\mathrm{GJ}}$ and $w_{4}(B - L)$ that the cited propositions state, computes the projection in (i) over all forty-five bivectors and checks $[\mu(T), T] = \tfrac58 T$, verifies (ii) for $p = \mu^{k}$, $k \le 5$, and verifies (iii) and (iv) as operator identities and entry by entry on the sixteen.
